\centerline{\bf \Huge Трапеция}

\begin{flushright}
        \scriptsize{}
\end{flushright}

\problem{1}
Углы при одном из оснований трапеции равны трапеции равны $50^{\circ}$ и $80^{\circ}$, а основания равны 2 и 3 . Найдите боковую сторону при угле $80^{\circ}$.

\problem{2}
Биссектриса одного угла трапеции делит её боковую сторону пополам. Найдите другую боковую сторону трапеции, если основания равны $a$ и $b$.

\problem{3}
Меньшая боковая сторона прямоугольной трапеции равна 1 , а большая образует угол $30^{\circ}$ с одним из оснований. Найдите это основание, если на нём лежит точка пересечения биссектрис углов при другом основании.

\problem{4}
Точка $M$ - середина стороны $C D$ параллелограмма $A B C D$. Точка $K$ на стороне $B C$ такова, что $\angle A M K=90^{\circ}$. Найдите $A K$, если $B K=a$ и $C K=b$.

\problem{5}
Боковая сторона трапеции равна одному основанию и вдвое меньше другого. Докажите, что вторая боковая сторона перпендикулярна одной из диагоналей.

\problem{6}
Трапеция $A B C D$ с основаниями $B C$ и $A D$ такова, что угол $A B D$ прямой и $B C+C D=A D$. Найдите отношение оснований $A D: B C$.